%%%%%%%%%%%%%%%%%%%%%%%%%%%%%%%%%%%%%%%%%
% Cover Letter - Progress in Artificial Intelligence
%%%%%%%%%%%%%%%%%%%%%%%%%%%%%%%%%%%%%%%%%

\documentclass{article}

\usepackage{charter}
\usepackage[
    a4paper,
    top=.5in,
    bottom=.5in,
    left=1in,
    right=1in,
]{geometry}

\setlength{\parindent}{0pt}
\setlength{\parskip}{0.45em}

\usepackage[colorlinks=true,linkcolor=magenta,citecolor=magenta,urlcolor=magenta]{hyperref}
\usepackage{fancyhdr}

\fancypagestyle{firstpage}{%
    \fancyhf{}
    \renewcommand{\headrulewidth}{0pt}
    \renewcommand{\footrulewidth}{1pt}
}

\AtBeginDocument{\thispagestyle{firstpage}}

\begin{document}

\vspace{-1em}
\rule{\linewidth}{1pt}

\bigskip

\hfill
\begin{tabular}{l @{}}
\today \bigskip\\
Mahmood Amintoosi\\
Computer Science Division\\
Department of Applied Mathematics \&\\
\href{https://gta-lab.github.io/}{Graph Theory and its Applications Lab}\\
Faculty of Mathematical Sciences,\\
Ferdowsi University of Mashhad,\\
IRAN,\\
Email: m.amintoosi@um.ac.ir\\
\url{https://mamintoosi.github.io/}
\end{tabular}

\bigskip

Dear Editor,

I am pleased to submit the manuscript entitled ``Graph Representation Learning for Student Sectioning: A Node2Vec-Based Approach'' for consideration as a research article in \textit{Progress in Artificial Intelligence}.

Student sectioning partitions enrolled students into instructional sections while preserving scheduling flexibility for subsequent course timetabling. Classical clustering approaches often represent students by binary enrollment vectors, which may miss relational structure from shared courses. We build a weighted student co-enrollment graph (after removing constant courses) and learn Node2Vec embeddings, then apply standard clustering without changing the downstream stage. The primary configuration ($p=1.0$, $q=0.5$, $d=2$) is selected by a quality--balance Pareto analysis validated over 20 seeds.

On six real-world courses (341 enrollments from 210 students), this setting reaches average Silhouette $0.619$, about $33\%$ above PCA+KMeans ($0.466$), the strongest baseline; the conventional representation ($0.157$) and spectral clustering ($0.128$) are weaker. Matched-$d$ PCA isolates graph learning from generic dimensionality reduction. Silhouette-optimal $d=1$ settings can collapse into imbalanced sections; the adopted setting improves quality and section balance together. The paper also reports grid search, sensitivity, stability, visualization, and runtime (under nine seconds per course). Limitations, including in-sample configuration selection on one institution, are stated.

The study is closely aligned with graph-learning topics previously published in the journal, including graph convolutional networks for clustering \cite{AlJreidy2024Clustering}, graph prototypical networks with active learning \cite{Solgi2022GraphPrototypical}, and link prediction based on local and global structure \cite{Luo2022LinkPrediction}. Our contribution complements these lines of work by applying random-walk-based graph representation learning to a classical educational clustering problem, with a controlled comparison against dimensionality reduction and direct graph clustering.

We believe the manuscript fits the journal's scope at the intersection of representation learning, clustering, graph-based machine learning, and AI applications in education. The manuscript is original, has not been published previously, and is not under consideration elsewhere. Source code, data, and experimental artifacts are available at \url{https://github.com/mamintoosi/Node2Vec-SSP}.

Thank you for considering this submission.

\bigskip

Sincerely,

Mahmood Amintoosi

\bibliographystyle{plain}
\bibliography{sn-bibliography}

\end{document}
